\documentclass[12pt,a4paper]{article}
\usepackage[utf8]{inputenc}
\usepackage[margin=1in]{geometry}
\usepackage{bookmark}
\usepackage{graphicx}
\usepackage{titlesec}
\usepackage{enumitem}
\usepackage{booktabs}
\usepackage{tabularx}
\usepackage{multirow}
\usepackage{tcolorbox}
\usepackage{setspace}
\usepackage{amsmath}
\usepackage{xcolor}
\usepackage{listingsutf8}

% ── Hyperlinks ────────────────────────────────────────────
\usepackage{hyperref}
\hypersetup{
    colorlinks=true,
    linkcolor=blue,
    filecolor=magenta,
    urlcolor=blue,
    pdftitle={Final Project Report: Intelligent Contract Risk Analysis},
}

% ── Listings (code snippets) ──────────────────────────────
\definecolor{codegray}{RGB}{247,247,247}
\definecolor{coderule}{RGB}{180,190,210}
\definecolor{codeblue}{RGB}{30,90,200}
\definecolor{codegreen}{RGB}{0,120,60}
\definecolor{codecomment}{RGB}{100,110,130}

\lstset{
    backgroundcolor=\color{codegray},
    basicstyle=\ttfamily\small,
    breaklines=true,
    keywordstyle=\color{codeblue}\bfseries,
    commentstyle=\color{codecomment}\itshape,
    stringstyle=\color{codegreen},
    numbers=none,
    frame=single,
    framesep=6pt,
    rulecolor=\color{coderule},
    xleftmargin=10pt,
    xrightmargin=10pt,
    captionpos=b,
    tabsize=4,
    columns=flexible,
    showstringspaces=false,
}

\setstretch{1.15}
\setlength{\parskip}{0.5em}

% ── Title ─────────────────────────────────────────────────
\title{
    \vspace{-1.5cm}
    \textbf{\Huge Intelligent Contract Risk Analysis} \\
    \vspace{0.6cm}
    \Large Final Report: Milestone 1 \& Milestone 2 \\
    \vspace{0.3cm}
    \large Milestone 1: ML-Based Contract Risk Classification \\
    \large Milestone 2: Agentic AI Legal Assistance System \\
    \vspace{0.4cm}
    \normalsize \textit{Submitted in fulfillment of the B.Tech CSE (AI) Capstone Project} \\
    \vspace{-0.2cm}
}

\author{
    \textbf{Newton School of Technology (NST), Sonipat} \\
    \vspace{0.8cm} \\
    \textbf{Team Members:} \\
    Ashutosh Singh (2401010109) \\
    Ranvendra Pratap Singh (2401010373) \\
    Shreya Suman (2401020068) \\
    Avishkar Meher (2401010116)
}
\date{April 2026}

% ─────────────────────────────────────────────────────────
\begin{document}

\maketitle
\thispagestyle{empty}
\vspace{1cm}

\begin{tcolorbox}[colback=gray!5, colframe=gray!50,
    title=\textbf{Declaration of Technical Integrity}]
\textbf{The ``No GenAI'' Rule:} We, the undersigned team members, formally affirm
that the core logic of this project is our original work. All data parsing modules,
Natural Language Processing pipelines, feature engineering implementations
(TF-IDF vectorization), Scikit-Learn model training, and the Streamlit inference
engine were manually and explicitly written and engineered by the team. The
LangGraph agentic pipeline (Milestone 2) was designed and implemented
independently, invoking external LLM APIs only for inference---not for generating
any project code or logic. We have adhered strictly to the grading rubric throughout
both milestones.
\end{tcolorbox}

\newpage
\tableofcontents
\newpage

% =============================================================
\section{Problem Statement}
\label{sec:problem}
% =============================================================

Modern corporate entities and legal firms process thousands of complex legal
contracts quarterly, ranging from Non-Disclosure Agreements (NDAs) to multi-million
dollar vendor agreements. The manual review of these documents is notoriously
labour-intensive, time-consuming, and highly susceptible to human error. Legal
professionals must meticulously scan dozens of pages of dense legal jargon to
locate and evaluate specific clauses related to indemnification, financial
liabilities, non-competes, and arbitrary termination conditions.

The objective of this project is to design, implement, and deploy an intelligent,
end-to-end contract analysis system that operates across two progressive milestones:

\begin{itemize}[leftmargin=1.5cm]
    \item \textbf{Milestone 1:} A supervised machine-learning pipeline that accepts
    raw unstructured contract text, intelligently segments it into individual clauses,
    and accurately classifies each clause into one of three actionable risk categories
    (\textbf{High, Medium, Low}).

    \item \textbf{Milestone 2:} An agentic AI system built on \textbf{LangGraph}
    that extends the ML pipeline with domain detection, Retrieval-Augmented Generation
    (RAG) from a legal knowledge base, and structured LLM-generated analysis reports
    including plain-English summaries, negotiation strategies, and safer clause
    rewrites.
\end{itemize}

By flagging high-severity clauses and providing quantitative confidence scores and
actionable legal guidance, this system drastically reduces the initial review burden
for legal teams.

% =============================================================
\section{Data Description}
\label{sec:data}
% =============================================================

\textbf{Source \& Provenance:} The machine learning pipeline was trained and
evaluated using the \textbf{Contract Understanding Atticus Dataset (CUAD) v1}, a
high-quality open-source corpus curated by expert legal professionals for NLP
research in the legal domain.

\vspace{0.2cm}
\noindent\textbf{Corpus Structure:} The raw dataset consists of a nested JSON
architecture containing \textbf{510 unique, fully annotated commercial contracts}.
Each document is structured hierarchically with 41 distinct legal clause types
annotated via expert Question-Answer (QA) pairs.

\vspace{0.2cm}
\noindent\textbf{Extraction Schema:} Our custom data pipeline parsed the JSON
natively. It filtered out all text blocks where the QA \texttt{is\_impossible} flag
was set to \texttt{true} (indicating that clause did not exist in the contract). For
valid clauses, the system extracted the \texttt{clause\_type} and the exact legal
text string from the \texttt{answers} key. Final processed output: \textbf{7.1 MB CSV}
containing 12,679 viable clause-level training examples.

\begin{table}[h]
\centering
\begin{tabular}{@{}ll@{}}
\toprule
\textbf{Property} & \textbf{Value} \\
\midrule
Total Contracts & 510 commercial contracts \\
Clause Types & 41 distinct legal clause categories \\
Source Format & JSON (structured annotation) \\
Domain Coverage & Commercial, Employment, NDA, SaaS, Vendor, Lease \\
Annotation Method & Expert legal professionals \\
Total Extracted Clauses & 12,679 (after filtering) \\
Processed CSV Size & 7.1 MB \\
\bottomrule
\end{tabular}
\caption{CUAD v1 Dataset Summary}
\end{table}

% =============================================================
\section{Exploratory Data Analysis}
\label{sec:eda}
% =============================================================

During the EDA phase, we systematically traversed the entire dataset to understand
the fundamental statistical properties of the legal text corpus. Several critical
insights shaped our methodology:

\begin{itemize}[leftmargin=1.5cm]
    \item \textbf{High Cardinality and Severe Class Imbalance:} The dataset natively
    contains 41 specific legal clause types. General boilerplate clauses (such as
    \textit{Governing Law}, \textit{Signatures}) appeared at drastically higher
    frequencies than severe penal clauses (like \textit{Liquidated Damages} or
    \textit{Uncapped Liability}).

    \item \textbf{Variable Context Lengths:} Significant variance was observed in
    text token lengths. Some clauses consisted of a single sentence, others (like IP
    Ownership clauses) spanned multiple dense paragraphs.

    \item \textbf{Contextual Ambiguity:} The same clause type could carry different
    effective risk levels depending on specific terms used (e.g., ``mutual'' vs.\
    ``unilateral'' indemnity).
\end{itemize}

\subsection*{Strategic Risk Mapping}

Predicting 41 heavily imbalanced classes using classical ML algorithms leads to
severe overfitting. To resolve this, we engineered a mapping architecture that
consolidated the 41 clause types into three actionable risk tranches:

\begin{enumerate}
    \item \textbf{High Risk:} Clauses inflicting severe financial, operational, or
    legal threats (e.g., Uncapped Liabilities, Termination for Convenience,
    Non-Competes, IP Ownership Assignments, Indemnification).

    \item \textbf{Medium Risk:} Clauses establishing operational friction or moderate
    constraints (e.g., Audit Rights, Insurance Requirements, Auto-Renewal,
    Arbitration clauses).

    \item \textbf{Low Risk:} Standard administrative or boilerplate language carrying
    no active threat (e.g., Governing Law, Dates, Notice Requirements).
\end{enumerate}

This consolidated the data into a healthy 3-class distribution:

\begin{table}[h]
\centering
\begin{tabular}{@{}lcc@{}}
\toprule
\textbf{Risk Class} & \textbf{Clause Count} & \textbf{Proportion} \\
\midrule
Low Risk & 5,629 & 44.4\% \\
Medium Risk & 4,340 & 34.2\% \\
High Risk & 2,710 & 21.4\% \\
\midrule
\textbf{Total} & \textbf{12,679} & \textbf{100\%} \\
\bottomrule
\end{tabular}
\caption{Risk Class Distribution After Consolidation}
\end{table}

% =============================================================
\section{Milestone 1: ML-Based Contract Risk Classification}
\label{sec:milestone1}
% =============================================================

\begin{tcolorbox}[colback=gray!5, colframe=gray!50]
\textbf{Objective:} Design and implement a supervised ML-based contract analysis
system that classifies clauses by risk level using classical NLP techniques,
\textit{without the use of Large Language Models or Generative AI}.
\end{tcolorbox}

\subsection{Preprocessing Pipeline}

Raw legal text is inherently noisy and full of formatting artifacts. Our cleaning
function executed the following systematic standardisations:

\begin{itemize}[leftmargin=1.5cm]
    \item \textbf{Case Normalisation:} All text transformed to lowercase for term
    uniformity.

    \item \textbf{Regex Scrubbing:} Stripped extraneous whitespace and special
    non-alphanumeric symbols, while intentionally retaining essential punctuation
    (\texttt{.,;:-'"()}) to preserve structural clause boundaries.

    \item \textbf{Clause Segmentation:} Contracts split on double-newline boundaries
    and sub-clause markers (numbered and lettered), retaining only segments
    exceeding 50 characters to filter out headings and formatting artifacts.
\end{itemize}

\subsection{Feature Engineering: TF-IDF Vectorisation}

To make text comprehensible to learning algorithms, we utilised \textbf{Term
Frequency-Inverse Document Frequency (TF-IDF)}:

\begin{equation}
\text{TF-IDF}(t, d, D) =
\underbrace{\frac{f_{t,d}}{\sum_{t' \in d} f_{t',d}}}_{\text{Term Frequency}}
\times
\underbrace{\log\frac{|D|}{|\{d \in D : t \in d\}|}}_{\text{Inverse Document Frequency}}
\end{equation}

\begin{enumerate}
    \item \textbf{Vocabulary Constraints:} \texttt{max\_features=5000}. This
    dimensionality reduction prevents the model from memorising rare names or dates
    unique to specific contracts, forcing it to learn actual legal syntax.

    \item \textbf{N-Gram Architecture:} \texttt{ngram\_range=(1,2)}. Captures
    standalone legal terms (unigrams: ``liability'') and essential contextual
    pairings (bigrams: ``breach of'', ``not exceed'', ``force majeure'').

    \item \textbf{Stop Word Removal:} Standard English stop words removed to
    eliminate structural noise.
\end{enumerate}

\subsection{Algorithms Evaluated}

\begin{enumerate}
    \item \textbf{Logistic Regression:} Selected for its efficiency in high-dimensional
    sparse spaces (5000-feature TF-IDF matrices). Provides calibrated probabilistic
    outputs via \texttt{predict\_proba}, enabling the custom confidence thresholding
    feature in the UI.

    \item \textbf{Decision Tree Classifier:} Evaluated to determine whether
    non-linear hierarchical branching rules outperformed linear boundaries.

    \item \textbf{Random Forest Classifier:} Bootstrap-aggregation ensemble specifically
    evaluated to combat the high variance and overfitting observed in isolated
    Decision Trees.
\end{enumerate}

\subsection{Optimisation Strategies}

\begin{enumerate}[label=\textbf{\arabic*.}]
    \item \textbf{Pipeline Architecture (Preventing Data Leakage):} The
    \texttt{TfidfVectorizer} was encapsulated directly \emph{inside} a Scikit-Learn
    \texttt{Pipeline} object alongside the classifier. This guarantees the vectorizer
    fits only on the training split of every cross-validation fold, never seeing
    holdout data.

    \item \textbf{Class Weighting:} Logistic Regression was instantiated with
    \texttt{class\_weight='balanced'}. This adjusts the internal cost function to
    penalise misclassifications on the minority High Risk class more severely than
    errors on the majority Low Risk class.

    \item \textbf{5-Fold Stratified Cross-Validation:} Stratification ensured the
    relative class distribution of High/Medium/Low clauses remained strictly identical
    across all training and holdout folds, producing statistically reliable generalisa\-tion estimates.
\end{enumerate}

\subsection{Evaluation Results}

\begin{table}[h]
\centering
\begin{tabular}{@{}lccc@{}}
\toprule
\textbf{Model Architecture} & \textbf{CV Accuracy} & \textbf{CV F1 (Weighted)} & \textbf{CV Precision} \\
\midrule
\textbf{Logistic Regression} & \textbf{88.56\%} & \textbf{0.8859} & \textbf{0.8861} \\
Decision Tree Classifier & 81.17\% & 0.8084 & 0.8091 \\
Random Forest Classifier & 81.38\% & 0.8058 & 0.8176 \\
\bottomrule
\end{tabular}
\caption{Model Comparison --- 5-Fold Stratified Cross-Validation}
\end{table}

\textbf{Why F1-Score?} While Accuracy is reported, we relied on the \textbf{Weighted
F1-Score} as the determinant metric. In legal risk analysis, False Negatives
(missing a High-Risk liability cap) are legally disastrous. The F1 harmonic mean
properly balances precision/recall tradeoffs.

\textbf{Final Selection:} Logistic Regression decisively outperformed tree-based
methods with a near 89\% F1-Score and was serialised as \texttt{best\_model.pkl}
for production inference.

\subsection{Milestone 1 Streamlit Dashboard}

The trained model was deployed as an interactive Streamlit dashboard featuring:

\begin{itemize}[leftmargin=1.5cm]
    \item Drag-and-drop PDF/TXT contract upload with instant text extraction via PyPDF2.
    \item Real-time metric cards showing High/Medium/Low clause counts and average
    confidence.
    \item Plotly Express visualisations: Risk distribution pie chart, confidence score
    histogram, risk intensity timeline, and top-keywords bar chart.
    \item Paginated, collapsible clause cards with risk labels, confidence percentages,
    and expandable full text.
    \item Confidence Threshold Slider: Dynamic filter to surface only
    high-confidence predictions (e.g., only flag clauses with $>$0.70 probability).
    \item One-click CSV export of the complete classified results.
\end{itemize}

% =============================================================
\section{Milestone 2: Agentic AI Legal Assistance System}
\label{sec:milestone2}
% =============================================================

\begin{tcolorbox}[colback=gray!5, colframe=gray!50]
\textbf{Objective:} Extend the ML classification system into a fully agentic AI
legal assistance tool that autonomously reasons about contract risks, retrieves
domain-specific legal best practices from a vector knowledge base, and generates
structured, actionable reports using Large Language Models.
\end{tcolorbox}

\subsection{Architecture Overview: From Static ML to Agentic AI}

Milestone 2 transforms the system from a \textbf{static ML classifier} into a
\textbf{fully agentic AI pipeline}. The trained Logistic Regression model from
Milestone 1 is now the \emph{entry point} of a multi-step autonomous reasoning
pipeline---not the final output.

The pipeline is orchestrated by a \textbf{LangGraph StateGraph} with four sequential
nodes sharing a single immutable \texttt{AgentState} TypedDict object:

\begin{center}
\begin{tabular}{@{}ccccc@{}}
\textbf{[S]} $\rightarrow$ &
\textbf{Node 1} & $\rightarrow$ & \textbf{Node 2} & $\rightarrow$ \\
& Domain Detect & & ML Classify & \\[4pt]
\textbf{Node 3} & $\rightarrow$ & \textbf{Node 4} & $\rightarrow$ & \textbf{[E]} \\
RAG Research & & LLM Reason & & Output \\
\end{tabular}
\end{center}

The key architectural property is \texttt{total=False} on the TypedDict, meaning
each node writes only the keys it produces, enabling clean error propagation and
full traceability at every stage. The compiled \texttt{StateGraph} is cached via
\texttt{@lru\_cache(maxsize=1)} to avoid re-compilation on every request.

\subsection{Node 1: Contract Domain Detection}

Before clause-level analysis begins, the system determines the contract's legal
domain. This label is used to filter RAG results, inject domain-specific context
into LLM prompts, and tag all outputs. Detection uses a two-stage algorithm:

\begin{enumerate}
    \item \textbf{Stage 1 --- Keyword Scoring (Zero API Cost):} A weighted regex
    signal table is applied against the first 3,000 characters of the document.
    Each of the 5 specific domains (NDA, Employment, Lease, SaaS, Vendor) has
    curated patterns with confidence weights on a 1--5 scale. Score $\geq 8$ = high
    confidence; $\geq 4$ = medium confidence.

    \item \textbf{Stage 2 --- LLM Fallback (When Score $<$ 4):} A single lightweight
    LLM call is made using Groq (free tier) with \texttt{max\_tokens=5} to
    disambiguate. The prompt instructs the model to return only one word from:
    \texttt{NDA | Employment | Lease | SaaS | Vendor | General}.
\end{enumerate}

Six domains are supported: NDA, Employment, Lease, SaaS, Vendor, and General Commercial.

\subsection{Node 2: ML Risk Classification}

This node bridges Milestone 1 and Milestone 2. It loads the pre-trained Logistic
Regression pipeline from \texttt{models/best\_model.pkl} and:

\begin{itemize}[leftmargin=1.5cm]
    \item Segments the contract via \texttt{segment\_clauses()}.
    \item Runs \texttt{model.predict()} and \texttt{model.predict\_proba()} to get
    risk label and confidence.
    \item Applies the user-specified confidence threshold.
    \item Forwards \textbf{only High and Medium} risk clauses to downstream nodes,
    skipping Low-risk clauses to conserve LLM API tokens.
    \item Truncates clauses exceeding 2,000 characters to stay within LLM context limits.
\end{itemize}

\subsection{Node 3: Domain-Aware RAG Research}

For each flagged clause, the Research node retrieves the top-3 most relevant legal
best-practice excerpts from the knowledge base. This retrieved context is injected
into the LLM prompt to ground its reasoning in established legal standards and
prevent hallucination.

The retrieval system uses a \textbf{3-Level Fallback Stack}:

\begin{enumerate}
    \item \textbf{Level 1 --- ChromaDB domain-filtered:} Highest quality. Queries
    the persistent vector store filtering by detected domain metadata. Embeddings
    produced by \texttt{all-MiniLM-L6-v2} (sentence-transformers; 384 dimensions).
    Distance metric: cosine similarity.

    \item \textbf{Level 2 --- ChromaDB full-collection:} If the domain has fewer
    results than the top-k target, the query is re-run across all domains.

    \item \textbf{Level 3 --- TF-IDF Fallback:} If ChromaDB is unavailable (e.g.,
    first cold start), an in-memory TF-IDF retriever over
    \texttt{data/legal\_best\_practices.json} serves results using cosine similarity
    with a topic-hint injection mechanism.
\end{enumerate}

The ChromaDB knowledge base contains \textbf{33+ domain-specific legal guidelines}
covering all 6 domains: NDA (6 entries), Employment (6), Lease (6), SaaS (6),
Vendor (4), and General Commercial (5).

\subsection{Node 4: LLM Deep Analysis (Reason Node)}

The Reason node dispatches each researched clause to a Large Language Model to
produce an \textbf{8-field structured analysis}. Every LLM response is parsed into
exactly the following fields:

\begin{table}[h]
\centering
\begin{tabularx}{\textwidth}{lX}
\toprule
\textbf{Field} & \textbf{Description} \\
\midrule
\texttt{plain\_english\_summary} & What the clause says at a 7th--8th grade reading level \\
\texttt{what\_makes\_it\_risky} & Specific legal danger with a real-world consequence \\
\texttt{who\_bears\_the\_risk} & Client / Vendor / Both / Third Party \\
\texttt{severity\_rationale} & Why the ML model rated this clause High or Medium \\
\texttt{industry\_standard\_practice} & What a fair, balanced version looks like \\
\texttt{negotiation\_tips} & 2--3 specific, actionable negotiation points \\
\texttt{safer\_rewrite} & Complete, ready-to-use formal replacement clause \\
\texttt{action\_required} & Remove / Negotiate / Seek Legal Review / Accept with Caution \\
\bottomrule
\end{tabularx}
\caption{LLM Analysis Output Schema --- 8 Structured Fields}
\end{table}

\textbf{Prompt Engineering:} The system prompt injects: (1) a ``Senior Legal Risk
Analyst AI'' persona, (2) the detected contract domain, (3) the ML risk label and
confidence score, (4) the retrieved RAG excerpts, and (5) strict anti-hallucination
rules forbidding citation invention. Temperature is set to 0.15 for consistent,
fact-grounded outputs. JSON output is enforced via
\texttt{response\_format=\{``type'': ``json\_object''\}}.

\textbf{Fault-Tolerant Parser (\texttt{safe\_parse\_analysis()}):} The JSON
parser is designed to never raise an exception. It handles:
\texttt{<think>...</think>} tags (Qwen/DeepSeek models), Markdown code fences
(\texttt{```json}), invalid JSON (regex substring search and retry), complete parse
failures (returns defaults for all 8 fields), and backward compatibility with
older 3-key response schemas.

\subsection{Cloud LLM --- 5-Level Resilient Fallback Chain}

The cloud LLM client implements a 5-level fallback chain ensuring the system
\emph{always} returns a result:

\begin{table}[h]
\centering
\begin{tabular}{@{}clll@{}}
\toprule
\textbf{\#} & \textbf{Provider} & \textbf{API Key Required?} & \textbf{Model} \\
\midrule
1 & \textbf{Pollinations AI} & No (free, no key) & \texttt{openai} \\
2 & OpenRouter (Key 1) & Yes & \texttt{openai/gpt-oss-120b} \\
3 & OpenRouter (Key 2) & Yes & \texttt{openai/gpt-oss-120b} \\
4 & DeepSeek via OpenRouter & Yes & \texttt{deepseek/deepseek-chat} \\
5 & Groq & Yes (free tier) & \texttt{llama-3.3-70b-versatile} \\
\midrule
$\times$ & \textit{All failed} & --- & \texttt{safe\_parse\_analysis} returns defaults \\
\bottomrule
\end{tabular}
\caption{Cloud LLM Fallback Chain --- Priority Order}
\end{table}

Pollinations AI (Level 1) requires \emph{no API key}, so the application works
end-to-end out of the box, satisfying the free-tier only constraint.

\subsection{Offline Mode --- Local Ollama LLM}

For complete privacy, the system supports a fully local offline mode powered by
Ollama:

\begin{itemize}[leftmargin=1.5cm]
    \item If the Ollama daemon is not running, the system auto-starts it via
    \texttt{subprocess.Popen(['ollama', 'serve'])}.
    \item Polls \texttt{GET /api/tags} every 500ms for up to 10 seconds to confirm
    readiness.
    \item Default model: \texttt{qwen3.5:2b} --- compact, efficient, excellent JSON output.
    \item Timeout: 180 seconds to accommodate slow first-load of quantised models on CPU.
\end{itemize}

\subsection{Report Generation}

Every completed analysis produces two types of output:

\begin{enumerate}
    \item \textbf{Structured JSON Report:} Machine-readable output containing the
    contract overview, detected domain, risk severity breakdown (High/Medium/Low
    counts), and per-clause analysis with all 8 LLM fields, ML metadata, and RAG
    references.

    \item \textbf{PDF Report (fpdf2):} A professionally formatted PDF with a branded
    header, per-clause colour-coded risk badges (red/amber/green), all 8 analysis
    fields with section headers, and a legal disclaimer on every page.
\end{enumerate}

\subsection{Milestone 2 Streamlit UI Enhancements}

The Milestone 2 UI adds a second tab (``AI Deep Analysis'') to the existing
Milestone 1 dashboard, featuring:

\begin{itemize}[leftmargin=1.5cm]
    \item Auto-detected domain badge with colour coding.
    \item Real-time clause-by-clause streaming progress indicator.
    \item Master-detail layout: clause list (left panel) and full analysis (right panel)
    with risk-severity coloured borders.
    \item JSON and PDF download buttons for the complete analysis report.
    \item Online/Offline toggle to switch between cloud LLM and local Ollama.
    \item RAG status pill showing ChromaDB (primary) or TF-IDF fallback (active retriever).
\end{itemize}

% =============================================================
\section{System Architecture \& Code Structure}
\label{sec:architecture}
% =============================================================

The repository is organised as a modular Python package:

\begin{lstlisting}
Contract-Risk-Classification/
|-- app.py                      # Main Streamlit application
|-- rag_setup.py                # ChromaDB vector store builder
|-- requirements.txt            # Python dependencies
|-- run_app.sh                  # Automated setup & launch script
|
|-- contract_agent/             # Core package (Milestone 2)
|   |-- workflow.py             # LangGraph StateGraph (4 nodes)
|   |-- core/
|   |   |-- state.py            # AgentState TypedDict
|   |   |-- domain.py           # Domain detection
|   |   |-- report.py           # JSON + Markdown reports
|   |   |-- pdf_report.py       # PDF export (fpdf2)
|   |-- llm/
|   |   |-- cloud.py            # 5-level LLM fallback chain
|   |   |-- local.py            # Ollama integration
|   |   |-- prompting.py        # Prompts + safe JSON parser
|   |-- retrieval/
|   |   |-- chroma.py           # ChromaDB retriever
|   |   |-- tfidf.py            # TF-IDF fallback retriever
|   |-- utils/
|       |-- ml.py               # Sklearn pipeline loader
|       |-- text.py             # clean_text, segment_clauses
|
|-- data/
|   |-- legal_best_practices.json   # Static KB for TF-IDF
|   |-- processed/processed_contracts.csv  # CUAD data (7.1 MB)
|
|-- models/
|   |-- best_model.pkl          # Trained LogReg pipeline (327 KB)
|
|-- src/                        # Jupyter notebooks (data pipeline)
    |-- 1_inspect.ipynb  ...  6_phase2_agent_demo.ipynb
\end{lstlisting}

% =============================================================
\section{Deployment}
\label{sec:deployment}
% =============================================================

The application is deployed on \textbf{Streamlit Community Cloud} at a publicly
accessible URL, satisfying the mandatory hosting requirement (no localhost demos).
Key deployment properties:

\begin{itemize}[leftmargin=1.5cm]
    \item Runs the complete ML Quick Screening pipeline without external dependencies.
    \item Automatically initialises the ChromaDB knowledge base on first run.
    \item Uses Pollinations AI free endpoint for LLM analysis with no API keys required.
\end{itemize}

\textbf{Local Setup:}

\begin{lstlisting}[language=bash]
git clone https://github.com/ashyou09/Contract-Risk-Classification.git
cd Contract-Risk-Classification
python -m venv .venv && source .venv/bin/activate
pip install -r requirements.txt
python rag_setup.py          # Build ChromaDB (first run only)
streamlit run app.py
\end{lstlisting}

% =============================================================
\section{Evaluation Against Grading Criteria}
\label{sec:grading}
% =============================================================

\subsection*{Milestone 1 Criteria (25\%)}

\begin{table}[h]
\centering
\begin{tabularx}{\textwidth}{lX}
\toprule
\textbf{Criterion} & \textbf{How Addressed} \\
\midrule
NLP \& ML techniques & TF-IDF (bigrams, 5000 features) + Logistic Regression. 5-fold stratified CV. F1 = 0.8859. \\
Preprocessing \& Feature Engineering & Custom clause segmentation, case normalisation, stop-word removal, bigram TF-IDF. \\
Evaluation metrics & Accuracy, Weighted F1, Precision, Recall, Confusion Matrix all computed. \\
UI Usability \& Code structure & Streamlit dashboard with 6+ charts, confidence threshold slider, CSV export. Fully modular codebase. \\
\bottomrule
\end{tabularx}
\caption{Milestone 1 Grading Criteria Mapping}
\end{table}

\subsection*{Milestone 2 Criteria (30\%)}

\begin{table}[h]
\centering
\begin{tabularx}{\textwidth}{lX}
\toprule
\textbf{Criterion} & \textbf{How Addressed} \\
\midrule
Agentic Reasoning \& Explainability & 8-field structured LLM output with plain-English summaries, safer rewrites, and negotiation tips. Domain-aware prompting. \\
RAG Integration \& State Management & ChromaDB with domain-filtered queries + TF-IDF fallback. Full \texttt{AgentState} TypedDict with error propagation across 4 nodes. \\
Ethical Awareness \& Responsible AI & Legal disclaimer on all outputs. Anti-hallucination prompt guardrails. ML confidence displayed. Offline mode for privacy. \\
Deployment & Live public URL on Streamlit Community Cloud. Works with zero API keys via Pollinations AI. \\
\bottomrule
\end{tabularx}
\caption{Milestone 2 Grading Criteria Mapping}
\end{table}

% =============================================================
\section{Ethical Considerations \& Responsible AI}
\label{sec:ethics}
% =============================================================

This system was built with responsible AI principles embedded by design:

\begin{enumerate}[leftmargin=1.5cm]
    \item \textbf{Legal Disclaimer:} Every output (Markdown report, JSON, PDF, and UI
    sidebar) displays: \textit{``This report is generated by an academic AI research
    prototype. It is NOT legal advice and may contain errors. Always consult a
    qualified legal professional before acting on any contract clause analysis.''}

    \item \textbf{Confidence Transparency:} ML confidence scores are displayed
    alongside every risk classification, enabling users to assess model certainty
    before acting.

    \item \textbf{Anti-Hallucination Prompting:} The LLM is explicitly instructed not
    to invent legal citations and to reference only provided RAG knowledge base
    excerpts.

    \item \textbf{Human Oversight:} The system presents analysis and evidence. All
    final legal decisions explicitly remain with the human user.

    \item \textbf{Privacy Preservation:} Offline mode (Ollama) enables complete
    on-device processing with no data leaving the user's machine.

    \item \textbf{No Over-Reliance Guardrail:} The LLM prompt explicitly forbids
    responses consisting solely of ``consult a lawyer'' without substantive analysis.
\end{enumerate}

% =============================================================
\section{Team Contribution}
\label{sec:team}
% =============================================================

The workload was systematically divided across both milestones:

\begin{description}[style=multiline, leftmargin=4cm]

    \item[\textbf{Ashutosh Singh} \\ (2401010109)]
    Developed the custom data preprocessing scripts responsible for parsing the CUAD
    JSON architecture. Implemented the core risk-label mapping logic (41 classes to
    3-tier system). Architected the base Scikit-Learn Pipeline preventing data
    leakage, trained and evaluated all ML models, and contributed to the overarching
    Streamlit UI dashboard structure.

    \vspace{0.2cm}
    \item[\textbf{Ranvendra Pratap} \\ \textbf{Singh} \\ (2401010373)]
    Implemented the TF-IDF feature engineering strategy. Co-defined the pipeline
    architecture. Designed and implemented the LangGraph StateGraph (Milestone 2),
    the ChromaDB RAG integration, the 5-level cloud LLM fallback chain, the Ollama
    local LLM integration, and the primary architect of the Streamlit UI for both
    milestones.

    \vspace{0.2cm}
    \item[\textbf{Shreya Suman} \\ (2401020068)]
    Led execution of the Scikit-Learn cross-validation pipeline. Performed Stratified
    K-Fold evaluation to establish champion model metrics. Led the report generation
    module (JSON, Markdown, PDF outputs) and contributed heavily to Streamlit UI
    front-end design and integration.

    \vspace{0.2cm}
    \item[\textbf{Avishkar Meher} \\ (2401010116)]
    Led quality assurance checks and end-to-end integration testing across both
    milestones. Performed deployment validation on Streamlit Community Cloud.
    Participated in code review, workflow refinement, and project documentation.

\end{description}

% =============================================================
\section*{References}
% =============================================================
\addcontentsline{toc}{section}{References}

\begin{enumerate}[leftmargin=1.5cm, label={[\arabic*]}]
    \item Hendrycks, D., Burns, C., Chen, A., \& Ball, S. (2021). \textit{CUAD: An
    Expert-Annotated NLP Dataset for Legal Contract Review}. arXiv:2103.06268.

    \item Pedregosa, F., et al. (2011). Scikit-learn: Machine Learning in Python.
    \textit{Journal of Machine Learning Research}, 12, 2825--2830.

    \item LangChain, Inc. (2024). LangGraph: Building Stateful Multi-Actor
    Applications with LLMs. \url{https://github.com/langchain-ai/langgraph}

    \item Reimers, N., \& Gurevych, I. (2019). Sentence-BERT: Sentence Embeddings
    using Siamese BERT-Networks. \textit{Proceedings of EMNLP 2019}.

    \item Chroma. (2024). ChromaDB: The AI-native open-source embedding database.
    \url{https://www.trychroma.com/}

    \item Lewis, P., et al. (2020). Retrieval-Augmented Generation for
    Knowledge-Intensive NLP Tasks. \textit{Advances in NeurIPS 2020}.

    \item Sparck Jones, K. (1972). A statistical interpretation of term specificity
    and its application in retrieval. \textit{Journal of Documentation}, 28(1), 11--21.

    \item Streamlit Inc. (2024). Streamlit --- The fastest way to build data apps.
    \url{https://streamlit.io/}
\end{enumerate}

\end{document}
